%%%PAGE 138 rot=0 legibility=good summary=卫星(轨道半径2R)被地球遮挡的追及角/遮挡角问题：先作切线、相对角速度、t_max；圆心角与挡光时间公式
%%%BLOCK id=138-01 ch=05 sec=遮光类·遮挡时间 kind=method alt=04 cont=none y=0.03-0.37
%FIG p138_a 0.09 0.04 0.31 0.20
\notefig[0.3]{figures/p138_a.png}

反向转

\[ \left(\frac{2\pi}{T}+\frac{2\pi}{T_0}\right)t_{\max}=\underline{\frac{\pi}{3}+\theta} \]
% 箭头从下面的“相对角速度.”指向上面括号内的 2π/T+2π/T_0

先作切线\quad 求追及角（遮挡角）
\[ (\omega_1+\omega_2)\,t=\underline{\unclear{\varphi}} \]% CHECK: 等号右侧为黑笔“α”形字符，上面又被红笔叠写成“φ”

相对角速度.

同向是 $\omega_1-(-\omega_2)=\omega_1+\omega_2$% CHECK: “同向/反向”与结果似乎对调（同向相对角速度应为 ω1-ω2），按原稿转录

反向是 $\omega_1-\omega_2=\omega_1-\omega_2$
%%%END
%%%BLOCK id=138-02 ch=05 sec=遮光类·遮挡时间 kind=method alt=04 cont=none y=0.42-0.78
%FIG p138_b 0.09 0.43 0.32 0.585
\notefig[0.3]{figures/p138_b.png}

\fbox{圆心角}\quad \red{$2\varphi=2\alpha+2\beta$}

挡光时间 \red{$t=\dfrac{2\varphi}{2\pi}\,T$}

\red{遮挡角 $=2\varphi$}
%%%END
%%%PAGEEND 138
